\section{Agent 失败模式与对策总结}

Anthropic 团队将所有观察到的失败模式及其对策整理为以下系统性对照表：

\begin{table}[H]
\centering
\small
\begin{tabular}{p{3cm}p{4.5cm}p{5.5cm}}
\toprule
\textbf{问题} & \textbf{Initializer Agent 行为} & \textbf{Coding Agent 行为} \\
\midrule
过早宣布项目完成 & 基于输入 spec 创建结构化的 JSON feature list & 每个 session 开始时阅读 feature list，选择单个功能开始工作 \\
\midrule
环境留有 bug 或未记录的进展 & 初始化 git 仓库和 progress notes 文件 & session 开始时阅读 progress 文件和 git log，运行基本测试捕捉 bug；session 结束时提交 git commit 并更新 progress \\
\midrule
过早标记功能为完成 & 创建包含验证步骤的 feature list & 所有功能都经过端到端测试后才标记为 passing \\
\midrule
花时间弄清楚如何运行项目 & 编写 \texttt{init.sh} 启动脚本 & session 开始时阅读 \texttt{init.sh} \\
\bottomrule
\end{tabular}
\caption{Agent 失败模式与双阶段对策的对应关系}
\end{table}

\begin{importantbox}{系统性思维：问题导向的设计}
这张对照表体现了一种优秀的工程实践——每个设计决策都直接对应一个具体的、已观察到的问题。这种"问题导向"的设计方法比"理论驱动"的方法更加实用，因为它确保每个机制都有明确的存在理由。在设计自己的 Agent harness 时，建议先系统性地收集和分类失败模式，再针对性地设计解决方案。
\end{importantbox}

\subsection{本章小结}

通过系统性地识别 Agent 的失败模式并分别在 Initializer Agent 和 Coding Agent 层面设计对策，Anthropic 团队构建了一个完整的问题-解决方案矩阵。每个设计决策都有明确的问题背景，这使得整个系统的设计理性且可解释。


